\section{Recomendaciones}
\begin{indentedsubsection}
Para asegurar el éxito a largo plazo de la plataforma Odoo y continuar con la evolución digital de Perú Flores Travel Agency, se proponen las siguientes recomendaciones operativas y estratégicas:
\end{indentedsubsection}

\subsection{Disciplina y exclusividad en el registro de información}
\begin{indentedsubsection}
Para que el sistema mantenga su valor y le ahorre tiempo real a la gerencia, es fundamental consolidar el hábito de registrar a los prospectos, armar las cotizaciones y anotar los adelantos de Western Union directamente en Odoo. Se recomienda abandonar de forma definitiva el uso paralelo de libretas físicas y archivos de Excel antiguos. La precisión y utilidad del sistema dependerá estrictamente de que no existan procesos comerciales ``ocultos'' fuera de la plataforma.
\end{indentedsubsection}

\subsection{Escalabilidad hacia la Integración Contable (SUNAT)}
\begin{indentedsubsection}
Si bien el proyecto actual ha delimitado claramente que la facturación electrónica continuará gestionándose de forma externa para no saturar la curva de aprendizaje inicial, se recomienda que, una vez la gerencia domine al 100\% el cotizador y el CRM, se evalúe en una segunda fase la integración directa con la SUNAT. Esto permitiría a futuro emitir las boletas y facturas con un solo clic desde la misma cotización ganada, cerrando el ciclo comercial y tributario en una sola pantalla.
\end{indentedsubsection}

\subsection{Aprovechamiento activo del ``Boca a Boca'' y Encuestas}
\begin{indentedsubsection}
Se sugiere establecer como regla obligatoria el envío del formulario de satisfacción (módulo de Encuestas) a cada turista extranjero inmediatamente después de finalizar su tour a destinos como Cusco, Arequipa o Puno. Recopilar estos testimonios de manera constante no solo servirá para medir la calidad de los guías y transportistas, sino que generará un historial de opiniones reales que la agencia podrá utilizar para capitalizar el ``boca a boca'' digital y atraer nuevos prospectos.
\end{indentedsubsection}

\subsection{Ciberseguridad y resguardo de datos del viajero}
\begin{indentedsubsection}
Al atender a turistas internacionales, la agencia maneja información sensible. Se recomienda aprovechar la seguridad de la nube de Odoo para almacenar cualquier documentación importante (como historial de ventas, preferencias de viaje o datos de contacto) dentro de la ficha de cada cliente en el CRM. Esto evitará que la información vital del negocio se quede vulnerable o se pierda si un teléfono celular se extravía o si se borran los chats de WhatsApp.
\end{indentedsubsection}
